\SourceLeaf{060}
\begin{center}
\pdfbookmark[2]{42.}{article-042}42.
\end{center}
We have constructed the table described above over a suitable range and added it to this work (Table I). The first part alone concerns the ellipse. We shall explain the second part, which covers hyperbolic motion, below. The argument of the table, the quantity $A$, runs from $0$ to $0{,}300$ in steps of one thousandth. The quantities $\log B$ and $C$ follow and are to be understood as expressed in $10\,000\,000^{\mathrm{ths}}$, or to seven decimal places, since the leading digits preceding the significant figures have been omitted. Finally, the fourth column shows the quantity $T$ computed first to 5, then to 6 figures. This precision is amply sufficient, since the column serves only to give, for the argument $T$, the corresponding values of $\log B$ and $C$ whenever we wish to determine $t$ from $v$ as in the preceding article. The converse problem, which occurs far more frequently, is to determine $v$ and $r$ from $t$ and can be solved entirely without the aid of $T$. We therefore preferred $A$ as the argument of our table to $T$, which would otherwise have been nearly as suitable and would even have made the table somewhat easier to construct. It is worth noting that all the numbers in the table were originally calculated to ten figures, so the seven figures given here can be relied on throughout. We cannot dwell here on the analytical methods used for this work, since a full explanation of them would take us too far from our purpose. The range of the table is amply sufficient for all cases in which it is advantageous to follow the method described so far. As shown above, beyond the limit $A=0{,}3$, which corresponds to $T=0{,}392374$ or $E=64^0\,8'$, one can conveniently dispense with special methods.
